In the CR5 pipeline, I aligned RGB timestamps, robot states, target actions, and gripper feedback before LeRobot conversion and policy serving. In the separate Hermite effort, I evaluated action-token representations under a controlled offline protocol. Across these independent efforts, I learned to trace how data definitions shape downstream evidence, but my experience remains with laboratory-scale and controlled pipelines. I now need systematic training in integration, governance, distributed scale, and data lineage. HKUST's MSc in Big Data Technology addresses that next step.

Provided the current offerings continue and MSBD6910 Independent Project receives approval, I would develop one lineage-centered workflow. MSBD5001 Foundations of Data Analytics would extend my collection, integration, cleansing, and schema checks toward governance, security, and privacy; MSBD5003 Big Data Computing would add distributed-computing and system-scale reasoning; and MSBD5012 Machine Learning would connect those data decisions to bounded model evidence. Within MSBD6910, I could design a scalable workflow that records timestamp provenance and representation transformations from synchronization through evaluation and serving. This would turn my existing pipeline experience into a clearer account of where data meaning changes, how each change is validated, and which evidence depends on it, preparing me to build dependable data infrastructure for embodied-agent R\&D.
